\documentclass{article}
\input{../../common_latex/import_include}
\input{../../common_latex/import_authors}
%-----------------------------------------------
% title
% 
\title{
	\Huge \textbf{NR Access} \\[1em]
}
\author{}
\date{}

%-----------------------------------------------
% document
%
\begin{document}

\maketitle

\tableofcontents
\newpage

\section{Terminology}
\input{../../common_latex/import_terms}

\section{介绍}
本节课程主要介绍NR接入流程。



\section{PRACH接入}
\subsection{工程经验}
UE侧的处理还会延迟约5个TA，物理层针对Preamble检测预留了2个TA的保护。如果考虑工程经验，覆盖范围还需要缩减。

\subsection{最大小区半径}
首先我们需要根据小区半径来选择PRACH Format来发送。
\begin{itemize}
\item Format 0$\to$3（仅支持FR1，详见TS38211 Table 6.3.3.2-2, Table 6.3.3.2-3, Table 6.3.3.2-4）\\
TS38211 5.3.2 OFDM baseband signal generation for PRACH原文如下"$N_{CP,l}^{RA}=N_{CP}^{RA}+n∙16k$ where for $\Delta f_{RA}={1.25,5}kHz$, n=0"\cite{38211}。这个说明对于format0$\to$3，原始的CP长度就是PRACH的CP长度。对于format0$\to$3，调度的slot是按照15kHz调度的\footnote{TS38211 6.3.3.2 Mapping to physical resources规定FR1按照15kHz计算slot数量，而format0$\to$3只占用FR1，且NR的调度单元是slot，所以一旦调度PRACH format 0$\to$3就会按照15kHz slot连续调度}。

为了避免PRACH自我干扰和干扰到其他调度单元，UE和BS空口传输时间$T_{prop}$必须满足
\begin{equation}
    min\{N_{CP}^{RA}, N_{seq}, N_{gap}\} \geq 2T_{prop}\footnotemark
\end{equation}
\footnotetext{这里是2$T_{prop}$：首先UE的时间窗口比BS的时间窗口晚了$T_{prop}$；为了克服UE发送中的时延，$T_{prop}$，UE相对自己的时间窗需要提前$T_{prop}$；需要相对BS的时间窗口提前2$T_{prop}$\label{fn:PRACH-2Tprop_reason}}
其中$N_{seq}$是PRACH中一个preamble的长度。

\begin{itemize}
\item Format 0: $N_{CP}^{RA} + N_u$=3168k+24576k=27744$k$=0.096875(ms), $N_{gap}$=1-0.096875=2976k\\
\[
r_{cell} = N_{gap}\cdot 3e8/2=14.5312(km)
\]
\item Format 1: $N_{CP}^{RA} + N_u$=21024k+2$\cdot$24576k=70176$k$=2.2844(ms), $N_{gap}$=3-2.2844=21984k
\[
r_{cell} = N_{CP}^{RA}\cdot 3e8/2=102.6562(km)
\]
\item Format 2: $N_{CP}^{RA} + N_u$=4688k+4$\cdot$24576k=102992$k$=3.3526(ms), $N_{gap}$=3.5-3.3526=4528k
\[
r_{cell} = N_{gap}\cdot 3e8/2=22.1094(km)
\]
\item Format 3: $N_{CP}^{RA} + N_u$=3168k+4$\cdot$24576k=27744$k$=0.096875(ms), $N_{gap}$=1-0.096875=2976k
\[
r_{cell} = N_{gap}\cdot 3e8/2=14.5312(km)
\]
\end{itemize}

\item Format A1$\to$A3支持\{15, 30, 60, 120, 480, 960\}kHz（没有$N_{gap}$）
\begin{itemize}
\item Format A1: $N_{CP}^{RA} + N_u$=$(288k+2\cdot 2048k)\cdot 2^{-\mu}$=$2\cdot(144k+2048k)\cdot 2^{-\mu}$, $N_{gap}$=0
\[
r_{cell} = N_{CP}^{RA}\cdot 3e8/2=\begin{cases}
1.4062(km) & scs=15kHz \\
0.7031(km) & scs=30kHz \\
0.3516(km) & scs=60kHz \\
0.1758(km) & scs=120kHz \\
0.0439(km) & scs=480kHz \\
0.0220(km) & scs=960kHz
\end{cases}
\]
\item Format A2: $N_{CP}^{RA} + N_u$=$(576k+4\cdot 2048k)\cdot 2^{-\mu}$=$4\cdot(144k+2048k)\cdot 2^{-\mu}$, $N_{gap}$=0
\[
r_{cell} = N_{CP}^{RA}\cdot 3e8/2=\begin{cases}
2.8125(km) & scs=15kHz \\
1.4062(km) & scs=30kHz \\
0.7031 (km) & scs=60kHz \\
0.3516(km) & scs=120kHz \\
0.0879(km) & scs=480kHz \\
0.0439(km) & scs=960kHz
\end{cases}
\]
\item Format A3: $N_{CP}^{RA} + N_u$=$(864k+6\cdot 2048k)\cdot 2^{-\mu}$=$6\cdot(144k+2048k)\cdot 2^{-\mu}$, $N_{gap}$=0
\[
r_{cell} = N_{CP}^{RA}\cdot 3e8/2=\begin{cases}
4.2188(km) & scs=15kHz \\
2.1094(km) & scs=30kHz \\
1.0547(km) & scs=60kHz \\
0.5273(km) & scs=120kHz \\
0.1318(km) & scs=480kHz \\
0.0659(km) & scs=960kHz
\end{cases}
\]
\end{itemize}

\item Format B1$\to$B4, Format C0,C2支持\{15, 30, 60, 120, 480, 960\}kHz（由于每次重复都占用一个OFDM symbol，那么所有slot中CP时间扣除$N_{CP}^{RA}$后就是$N_{gap}$）
\begin{itemize}
\item Format B1, $N_{CP}^{RA}$=$216k\cdot 2^{-\mu}$, $N_{gap}$=$72k\cdot 2^{-\mu}$
\[
r_{cell} = N_{gap}\cdot 3e8/2=\begin{cases}
0.3516(km) & scs=15kHz \\
0.1758(km) & scs=30kHz \\
0.0879(km) & scs=60kHz \\
0.0439(km) & scs=120kHz \\
0.0110(km) & scs=480kHz \\
0.0055(km) & scs=960kHz
\end{cases}
\]
\item Format B2, $N_{CP}^{RA}$=$360k\cdot 2^{-\mu}$, $N_{gap}$=$216k\cdot 2^{-\mu}$
\[
r_{cell} = N_{gap}\cdot 3e8/2=\begin{cases}
1.0547(km) & scs=15kHz \\
0.5273(km) & scs=30kHz \\
0.2637(km) & scs=60kHz \\
0.1318(km) & scs=120kHz \\
0.0330(km) & scs=480kHz \\
0.0165(km) & scs=960kHz
\end{cases}
\]
\item Format B3, $N_{CP}^{RA}$=$504k\cdot 2^{-\mu}$, $N_{gap}$=$360k\cdot 2^{-\mu}$
\[
r_{cell} = N_{gap}\cdot 3e8/2=\begin{cases}
1.7578(km) & scs=15kHz \\
0.8789(km) & scs=30kHz \\
0.4395(km) & scs=60kHz \\
0.2197(km) & scs=120kHz \\
0.0549(km) & scs=480kHz \\
0.0275(km) & scs=960kHz
\end{cases}
\]
\item Format B4, $N_{CP}^{RA}$=$936k\cdot 2^{-\mu}$, $N_{gap}$=$792k\cdot 2^{-\mu}$
\[
r_{cell} = N_{gap}\cdot 3e8/2=\begin{cases}
3.8672(km) & scs=15kHz \\
1.9336(km) & scs=30kHz \\
0.9668(km) & scs=60kHz \\
0.4834(km) & scs=120kHz \\
0.1208(km) & scs=480kHz \\
0.0604(km) & scs=960kHz
\end{cases}
\]
\item Format C0, $N_{CP}^{RA}$=$1240k\cdot 2^{-\mu}$, $N_{gap}$=$1096k\cdot 2^{-\mu}$
\[
r_{cell} = N_{CP}^{RA}\cdot 3e8/2=\begin{cases}
5.3516(km) & scs=15kHz \\
2.6758(km) & scs=30kHz \\
1.3379(km) & scs=60kHz \\
0.6689(km) & scs=120kHz \\
0.1672(km) & scs=480kHz \\
0.0836(km) & scs=960kHz
\end{cases}
\]
\item Format C2, $N_{CP}^{RA}$=$2048k\cdot 2^{-\mu}$, $N_{gap}$=$2912k\cdot 2^{-\mu}$
\[
r_{cell} = N_{CP}^{RA}\cdot 3e8/2=\begin{cases}
10(km) & scs=15kHz \\
5(km) & scs=30kHz \\
2.5(km) & scs=60kHz \\
1.25(km) & scs=120kHz \\
0.3125(km) & scs=480kHz \\
0.1562(km) & scs=960kHz
\end{cases}
\]
\end{itemize}
\end{itemize}

\subsubsection{计算当前小区适配的$N_{cs}$}
上面给出的最大小区半径是$N_{cs}=0$，如果小区半径在最大半径以内，则可以计算$N_{cs}$最小值以满足最大半径覆盖范围(如果考虑定时误差、多径时延冗余、安全冗余还需要增加$T_{sec}$)。

\begin{figure}[htbp]
    \centering
    \includegraphics[width=0.5\textwidth, keepaspectratio]{lect001_access_tdd_guard}
    \caption{TDD Guad时间}
    \label{fig:example}
\end{figure}
\[
N_{cs} = L_{RA}\lfloor \frac{2r_{cell}/c+T_{sec}}{T_{seq}} \rfloor
\]


\subsection{小区半径对TDD Slot结构的限制}
小区半径此时已经确认，TDD传输时DL$\to$UL切换时，需要预留保护间隔（GAP），确保上下行隔离，不存在交叠。如下图所示，GAP的长度要求大于下行延时量和上行提前量之和，即$2T_{prop}+N_{TA,offset}$，其中$T_{prop}$表示小区边缘用户的信号接收时延。
\[
T_{guard} \geq N_{TA,offset} + 2T_{prop}\textsuperscript{\ref{fn:PRACH-2Tprop_reason}}
\]
其中$T_{prop}$就是小区半径导致的传输延时。
\begin{longtable}{|p{0.15\textwidth}|p{0.15\textwidth}|p{0.15\textwidth}|}
\caption{PRACH Long Format $T_{guard} $}\\
\hline
 & \multicolumn{1}{|c|}{$N_{TA,offset}=25600$} & \multicolumn{1}{|c|}{$N_{TA,offset}=39936$} \\ \hline
Format 0 & 3376k & 3600k \\ \hline
Format 1 & 22384k & 22608k \\ \hline
Format 2 & 4928k & 5152k \\ \hline
Format 3 & 3376k & 3600k \\ \hline
\end{longtable}

\begin{longtable}{|p{0.15\textwidth}|p{0.15\textwidth}|p{0.15\textwidth}|p{0.15\textwidth}|}
\caption{PRACH Short Format $T_{guard} $}\\
\hline
& \multicolumn{1}{|c|}{$N_{TA,offset}=25600$} & \multicolumn{1}{|c|}{$N_{TA,offset}=39936$} & \multicolumn{1}{|c|}{$N_{TA,offset}=13792$} \\ \hline
Format A1 & 288k$\cdot 2^{-\mu}$+400 & 288k$\cdot 2^{-\mu}$+624 & 288k$\cdot 2^{-\mu}$+215.5 \\ \hline
Format A2 & 576k$\cdot 2^{-\mu}$+400 & 576k$\cdot 2^{-\mu}$+624 & 576k$\cdot 2^{-\mu}$+215.5 \\ \hline
Format A3 & 864k$\cdot 2^{-\mu}$+400 & 864k$\cdot 2^{-\mu}$+624 & 864k$\cdot 2^{-\mu}$+215.5 \\ \hline
Format B1 & 72k$\cdot 2^{-\mu}$+400 & 72k$\cdot 2^{-\mu}$+624 & 72k$\cdot 2^{-\mu}$+215.5 \\ \hline
Format B2 & 216k$\cdot 2^{-\mu}$+400 & 216k$\cdot 2^{-\mu}$+624 & 216k$\cdot 2^{-\mu}$+215.5 \\ \hline
Format B3 & 360k$\cdot 2^{-\mu}$+400 & 360k$\cdot 2^{-\mu}$+624 & 360k$\cdot 2^{-\mu}$+215.5 \\ \hline
Format B4 & 792k$\cdot 2^{-\mu}$+400 & 792k$\cdot 2^{-\mu}$+624 & 792k$\cdot 2^{-\mu}$+215.5 \\ \hline
Format C0 & 1096k$\cdot 2^{-\mu}$+400 & 1096k$\cdot 2^{-\mu}$+624 & 1096k$\cdot 2^{-\mu}$+215.5 \\ \hline
Format C2 & 2048k$\cdot 2^{-\mu}$+400 & 2048k$\cdot 2^{-\mu}$+624 & 2048k$\cdot 2^{-\mu}$+215.5 \\ \hline
\end{longtable}



\paragraph{$N_{TA,offset}$}\mbox{}\\
\begin{itemize}
\item $N_{TA,offset}$由高层RRC n-TimingAdvanceOffset（出现在ServingCellConfigCommon和\\
ServingCellConfigCommonSIB 中），有$n0, n25600, n39936$三个值（表示对应的数量的tc）\cite{38331}。
\item 高层RRC n-TimingAdvanceOffset没有配置时，UE采用38133给的默认值（TS38213 4.2Transmission timing adjustments的原文是"If the UE is not provided n-TimingAdvanceOffset for a serving cell, the UE determines a default value NTA,offset of the timing advance offset for the serving cell as described in [10, TS 38.133]"）\cite{38213}。
\begin{itemize}
\item TS38133 7.1.1	Introduction原文"$N_{TA}$ offset for first TAG and second TAG are as defined in n-TimingAdvanceOffset and n-TimingAdvanceOffset2 respectively, if configured, otherwise it is $N_{TA,offset}$ defined in table 7.1.2-2."\cite{38133}。这说明缺省时，采用Table 7.1.2-2配置。

\begin{longtable}{|p{0.8\textwidth}|p{0.19\textwidth}|}
\caption{TS38133 Table 7.1.2-2\cite{38133}}\\
\hline
\multicolumn{1}{|c|}{Frequency range and band of cell used for uplink transmission} & \multicolumn{1}{c|}{$N_{TA,offset}$(Unit: $T_c$)} \\ \hline
FR1 FDD or TDD band with neither E-UTRA–NR nor NB-IoT–NR coexistence case & 25600 \\ \hline
FR1 FDD band with E-UTRA–NR and/or NB-IoT–NR coexistence case & 0 \\ \hline
FR1 TDD band with E-UTRA–NR and/or NB-IoT–NR coexistence case & 39936 \\ \hline
FR2 & 13792 \\ \hline
\end{longtable}

\end{itemize}

\end{itemize}




\section{接入故障排查}
接入排查分为DL/UL RCT
\subsection{RF}
首先要进行RF测试：用Vector Signal Generator和BS Rx连接，测试接收机灵敏度在dBm下正确率有多少(e.g., -100dBm)。具体详见TS38104
\subsection{End to End}
如Fig~\ref{fig:accesss_process}所示，在接入流程中有竞争接入和非竞争接入
\begin{figure}[htbp]
    \centering
    \includegraphics[width=\textwidth, keepaspectratio]{lect001_access}
    \caption{Access Process}
    \label{fig:accesss_process}
\end{figure}

\subsubsection{竞争接入}
\paragraph{Power Control Issue}\mbox{}
\begin{itemize}
\item Msg1 Tx power 计算公式如下
\begin{align}
P_{loss} &= P_{rx}(SSB)-P_{tx}(SSB)\\ \notag
P_{tx}(Msg1) &= P_{rx}(Msg1) - P_{loss} \notag
\end{align}
$P_{tx}(SSB)$由SIB1中的ss-PBCH-BlockPower给出，$P_{rx}(Msg1)$由SIB1中preambleReceivedTargetPower给出。\\
{\colorbox{yellow}{提高覆盖距离的办法就是提高SIB1中的 preambleReceivedTargetPower}}
\item Msg3的计算公式如下
\[
P_{tx}(Msg3)=P_{tx}(Msg1) + TPC(Msg2)
\]
{\colorbox{yellow}{Msg2中的UL-GRANT的TPC降低发射功率会导致Msg3接收不到}}
\end{itemize}
\paragraph{TA测量问题}\mbox{}
\begin{itemize}
\item Msg1(PRACH)常见问题是preamble id error\\
对于SDR TDD平台，这个问题是Tx/Rx之间有时间差($N_{TA, offset}$=0, 25600$T_c$, 39936$T_c$, 13792$T_c$)，这个时间差导致超过$N_{cs}$ 窗口。\\
解决方法如下：
\begin{itemize}
\item set $N_{cs}=0$, 去测量偏移量，调整SDR/PC802里面对应的偏移量参数来对齐$N_{TA, offset}$
\item TA在测量时会偏移2~4个采样点（这是过采样后的sinc波导致的）{\colorbox{yellow}{Blackman窗或许可以解决此问题}}
\end{itemize}

\item Msg2(RAR)中会发送preamble id，UE需要找到自己的；TA用于发送Msg3；UL-Grant告诉Msg3的位置

\end{itemize}





\newpage
\bibliographystyle{IEEEtran}
\bibliography{../../common_latex/import_reference}
\end{document}
